\documentclass[conference]{IEEEtran}
\IEEEoverridecommandlockouts
\usepackage{cite}
\usepackage{amsmath,amssymb,amsfonts}
\usepackage{algorithmic}
\usepackage{graphicx}
\usepackage{textcomp}
\usepackage{xcolor}
\usepackage[ruled,vlined]{algorithm2e}
\usepackage{url}
\usepackage[utf8]{inputenc}
\usepackage{array}
\usepackage{float}
\usepackage{booktabs}
\usepackage{hyperref}
\usepackage{xurl}
\hypersetup{breaklinks=true}

\def\BibTeX{{\rm B\kern-.05em{\sc i\kern-.025em b}\kern-.08em
    T\kern-.1667em\lower.7ex\hbox{E}\kern-.125emX}}

\begin{document}

\title{Multimodal RAG for Disaster Intelligence: A Simulation-Driven Approach}

\author{
\IEEEauthorblockN{
Christopher Aseer J Albert\textsuperscript{1},
Ashutosh Monish Bhonsle\textsuperscript{2},
Prisha Andrew Bhosale\textsuperscript{3},
Harshvardhan Bhota\textsuperscript{4},
Chaitanya Vijaykumar Mahamuni\textsuperscript{5}
}
\IEEEauthorblockA{
\textsuperscript{1,2,3,4}\textit{Department of Computer Engineering,
Fr. C Rodrigues Institute of Technology, Vashi, India}\\
\textsuperscript{5}\textit{Electronics and Telecommunication Engineering,
University of Mumbai, Mumbai, India}\\
\textsuperscript{1}\texttt{christpheraseer5080@gmail.com},
\textsuperscript{2}\texttt{bhonsle.ashutosh@gmail.com},
\textsuperscript{3}\texttt{prisha.bhosale22@gmail.com},
\textsuperscript{4}\texttt{harshvardhanbohta@gmail.com},
\textsuperscript{5}\texttt{chaitanyamahamuni91@gmail.com}
}
}

\maketitle

\begin{abstract}
Effective disaster response depends on rapid, accurate situational awareness
derived from heterogeneous data sources including satellite imagery, social
media reports, environmental sensor readings, and official emergency
management documents. Existing systems typically process these modalities in
isolation, limiting their capacity for holistic reasoning in safety-critical
conditions. This paper presents a unified Multimodal Retrieval-Augmented
Generation (RAG) framework for disaster intelligence. The system integrates
an autonomous LangGraph ReAct agent with live hydrological and meteorological
APIs, a local vision-language model for satellite imagery analysis, and a
multi-modal vector knowledge base covering social reports, IoT telemetry, and
official protocols. A deterministic safety validator is incorporated as a
post-processing layer to detect and flag unsafe or ungrounded outputs. The
framework is evaluated on a 50-query benchmark drawn from real-world disaster
corpora across five scenario types. The proposed system achieves a Keyword F1
score of 0.336, DR-Fact of 0.624, and Situational Awareness Accuracy of 0.698,
outperforming both a knowledge-only LLM baseline and a text-only RAG
configuration on all primary metrics.
\end{abstract}

\begin{IEEEkeywords}
Disaster Response; Multimodal Learning; Retrieval-Augmented Generation;
Situational Awareness; Large Language Models; Safety-Critical AI
\end{IEEEkeywords}

% ── Section I ─────────────────────────────────────────────────────────────────
\section{Introduction}

Natural disasters including floods, earthquakes, wildfires, and hurricanes
cause substantial loss of life and infrastructure damage worldwide~\cite{b1}.
Effective response requires timely situational awareness to support emergency
decision-making and search-and-rescue operations~\cite{b2}. Modern disaster
environments produce large volumes of heterogeneous data from satellite and UAV
imagery, social media, environmental sensors, and official management
documents. Synthesising these disparate streams into actionable intelligence
remains a significant open problem~\cite{b3}.

Recent advances in artificial intelligence have improved specific aspects of
disaster monitoring, such as damage assessment from satellite imagery and
crisis-related text classification~\cite{b4}. However, most existing methods
address individual modalities or rely on late-fusion pipelines that perform
limited cross-modal reasoning. Large Language Models (LLMs) offer strong
generative and reasoning capabilities but are susceptible to hallucinations
when deployed without grounding mechanisms, particularly in high-stakes
domains~\cite{b5}. Retrieval-Augmented Generation (RAG) addresses this by
anchoring responses to retrieved evidence~\cite{b6}, yet existing RAG systems
are predominantly text-based and have not been applied systematically to
multimodal disaster intelligence.

This paper makes the following contributions:
\begin{enumerate}
  \item A unified multimodal RAG framework that integrates visual,
        textual, sensor, and protocol knowledge through an autonomous
        LangGraph ReAct agent for end-to-end disaster reasoning.
  \item A modality-aware retrieval and evidence-fusion strategy employing
        Maximal Marginal Relevance (MMR) across four indexed data types,
        combined with a Temporal Context Bridge to prevent conflation of
        historical and live information.
  \item A deterministic, regex-based safety validator that scans every
        generated response for unsafe phrases and failure-to-warn patterns
        before output is presented to the operator.
\end{enumerate}

The remainder of this paper is organised as follows. Section~II surveys
related work. Section~III describes the system design. Section~IV presents
the experimental evaluation. Section~V concludes the paper.

% ── Section II ────────────────────────────────────────────────────────────────
\section{Related Work}

\subsection{Visual and Multimodal Situational Awareness}
Satellite and UAV imagery are central to post-disaster damage assessment.
Singh and Hoskere~\cite{b15} demonstrated that transformer-based vision models
achieve state-of-the-art results on ultra-high-resolution aerial datasets for
damage classification. Ofli et al.~\cite{b16} showed that multimodal deep
learning, combining image and text from social media, outperforms unimodal
approaches for disaster event categorisation. However, most such systems
operate under closed-world assumptions and use late-fusion architectures that
limit cross-modal reasoning and adaptability to novel scenarios.

\subsection{Social Media and Sensor Data Integration}
Social media provides real-time situational intelligence during disasters but
is noisy and unstructured. Alam et al.~\cite{b18} introduced the HumAID
dataset and benchmarked deep learning classifiers for humanitarian tweet
categorisation. Abavisani et al.~\cite{b19} demonstrated that cross-modal
representation learning improves crisis event classification from multimodal
social media. Ray et al.~\cite{b20} highlighted the growing role of IoT-based
sensor networks in real-time disaster monitoring and decision support. Combining
social media analytics with physical sensor readings improves the credibility
of situational reports and reduces information overload for responders.

\subsection{Knowledge Graph and RAG Approaches}
Zhang et al.~\cite{b13} proposed combining ontology-driven knowledge graph
construction with RAG to support earthquake emergency response, integrating
authoritative documents through hybrid retrieval for improved factual accuracy.
Lewis et al.~\cite{b6} established RAG as a general paradigm for grounding
LLM outputs in retrieved external evidence for knowledge-intensive tasks.
While effective for text-based decision support, these methods do not
incorporate live sensor streams or visual modalities.

\subsection{Reliability and Safety in Disaster AI}
Cantini et al.~\cite{b10} studied the use of prompt-based LLMs for automated
disaster reporting from social media, noting reliability limitations from
hallucinations. Yu et al.~\cite{b14} demonstrated that retrieval-augmented and
hallucination-aware correction mechanisms improve accuracy in vision-language
spatial description tasks. Kim et al.~\cite{b12} proposed a multimodal
framework for real-time disaster event classification and prioritisation.
These works collectively motivate the need for grounding, safety verification,
and multi-source fusion in disaster-oriented AI systems.

% ── Section III ───────────────────────────────────────────────────────────────
\section{System Design}

\subsection{Architecture Overview}

The proposed system is built around a ReAct (Reason + Act) loop implemented
with LangGraph~\cite{b23}, powered by the Llama-3.3-70B language
model~\cite{b24} accessed via the Groq API. At each turn, the agent selects
and executes one or more specialised tools, incorporates the returned evidence
into its reasoning context, and generates a grounded response. Conversation
history is maintained across sessions using LangGraph's \texttt{MemorySaver}
checkpointing backend, with each session assigned an independent
\texttt{thread\_id} to prevent state leakage. The agent is configured with
temperature $T{=}0.1$ to favour deterministic and factually grounded outputs.

\subsection{Knowledge Base Construction}

Four data modalities are indexed into a single ChromaDB vector store.
Document chunks are embedded using the \texttt{BAAI/bge-small-en-v1.5}
model, which produces 384-dimensional L2-normalised vectors. For normalised
vectors, cosine similarity reduces to a dot product:

\begin{equation}
  S(q,d) = \sum_{i=1}^{384} E(q)_i \cdot E(d)_i
  \label{eq:cosine}
\end{equation}

Official protocol documents (FEMA and CDC manuals) are segmented using a
recursive character text splitter with chunk size $C_{\text{size}}{=}1{,}000$
characters and overlap $C_{\text{overlap}}{=}200$ characters, preserving
semantic continuity within the model's 512-token context limit.
Table~\ref{tab:kb} summarises the knowledge base composition.

\begin{table}[ht]
\centering
\caption{Knowledge base composition.}
\label{tab:kb}
\renewcommand{\arraystretch}{1.2}
\begin{tabular}{@{}llr@{}}
\toprule
\textbf{Modality Tag} & \textbf{Source} & \textbf{Documents} \\ \midrule
\texttt{text}                & CrisisLex-T26 (2013 CO Floods) & $\leq$500 \\
\texttt{imagery\_description}& xView2 + LLaVA-7B              & 50 \\
\texttt{telemetry}           & USGS NWIS Site 06730200        & $\approx$34 \\
\texttt{knowledge\_base}     & FEMA + CDC PDFs                & $\approx$3{,}000--5{,}000 \\ \bottomrule
\end{tabular}
\end{table}

Social media tweets were drawn from the CrisisLex-T26
dataset~\cite{b21}, filtered to the 2013 Colorado Floods subset. Tweet
timestamps were reconstructed from Twitter Snowflake identifiers to enable
temporal tagging. Satellite imagery was sourced from xView2~\cite{b22} and
processed locally by LLaVA-7B running under Ollama, producing structured
textual descriptions (SUMMARY, DAMAGE, ACCESS, FULL\_DESC) for each of 50
images. Live hydrological telemetry was retrieved from the USGS National Water
Information System (NWIS) API for gauge site 06730200 (Boulder Creek, Colorado),
with readings classified as NORMAL, WARNING, DANGER, or CRITICAL against the
site's 10.0\,ft flood-stage threshold.

\subsection{Retrieval Strategy}

At inference time, the agent issues semantic queries against the vector store
using Maximal Marginal Relevance (MMR) retrieval with fetch window
$k_f{=}20$, return set $k{=}7$, and diversity weight $\lambda{=}0.7$.
MMR balances relevance with intra-result diversity, reducing redundancy in
retrieved crisis documents~\cite{b25}. Each retrieved document carries a
metadata tag (\texttt{modality}, \texttt{source}, \texttt{timestamp},
\texttt{location}) enforced through ChromaDB's JSON schema contract. A
modality-filtering policy prevents operational FEMA guidance from being
conflated with social media reports during retrieval.

\subsection{Temporal Context Bridge}

A key architectural challenge is preventing the agent from treating historical
disaster records as current operational intelligence. Retrieved documents are
tagged with a temporal label $T_{\text{tag}}$ derived from their stored
timestamp metadata and concatenated to the document content before injection
into the context window:

\begin{equation}
  C(d) = T_{\text{tag}}(t_d) \oplus d
  \label{eq:tcb}
\end{equation}

Documents from the 2013 Colorado Floods corpus receive the tag
\texttt{[HISTORICAL -- 2013]}; live API outputs receive \texttt{[LIVE --
current]}. This mechanism is enforced at the retrieval layer without requiring
metadata-based timestamp filtering, which would suppress relevant historical
heuristics. The agent system prompt explicitly instructs the model to
prioritise live data for operational decisions and use historical data only
for pattern recognition.

\subsection{Tool Repertoire and Agentic Loop}

The agent selects from five specialised tools at each reasoning step.
Table~\ref{tab:tools} lists the tools and their functions.

\begin{table}[ht]
\centering
\caption{Agent tool repertoire.}
\label{tab:tools}
\renewcommand{\arraystretch}{1.2}
\begin{tabular}{@{}ll@{}}
\toprule
\textbf{Tool} & \textbf{Function} \\ \midrule
\texttt{fetch\_live\_weather}     & Open-Meteo real-time weather retrieval \\
\texttt{query\_water\_sensor}     & USGS NWIS live river gauge heights \\
\texttt{query\_recent\_earthquakes}& USGS seismic events within 300\,km \\
\texttt{search\_official\_protocols}& MMR search of FEMA/CDC knowledge base \\
\texttt{search\_social\_reports}  & MMR search of social and imagery KB \\
\texttt{analyze\_disaster\_image} & On-device LLaVA-7B image analysis \\ \bottomrule
\end{tabular}
\end{table}

Algorithm~\ref{alg:fallback} describes the agentic execution loop with
fault-tolerance handling. If a live API returns an error, the agent falls back
to a pre-configured alternative location or to historical vector-store
retrieval, and notifies the user of the fallback event.

\begin{algorithm}[ht]
\caption{Agentic API Execution with Fault Tolerance}
\label{alg:fallback}
\SetAlgoLined
\KwIn{Query $Q$, Location $L$}
\KwOut{Operational Response $R$}
Initialise ReAct state with $Q$\;
\eIf{$Q$ requires live telemetry}{
  Attempt \texttt{API\_Call}($L$)\;
  \uIf{success}{
    $\text{Context} \leftarrow \text{Parse}(\texttt{API\_Call})$\;
  }
  \uElseIf{error = ``Location Not Found''}{
    $\text{Context} \leftarrow \texttt{API\_Call}(\text{FallbackLocation})$\;
    Notify user of fallback\;
  }
  \Else{
    $\text{Context} \leftarrow \text{VectorSearch}(Q, \text{KB})$\;
    Notify user of historical fallback\;
  }
}{
  $\text{Context} \leftarrow \text{VectorSearch}(Q, \text{DB})$\;
}
$R \leftarrow \text{LLM\_Synthesise}(\text{Context}, Q)$\;
$R \leftarrow \text{SafetyValidator}(R)$\;
\Return{$R$}\;
\end{algorithm}

\subsection{Safety Validator}

All generated responses pass through a deterministic post-processing layer
before being returned to the operator. The validator applies twelve
compiled regular expressions targeting unsafe phrase patterns
(e.g., \textit{safe to drive through}, \textit{floodwater is clean},
\textit{no risk}) and separately checks responses to queries annotated as
dangerous for the presence of at least one explicit safety warning keyword.
Any flagged response is prepended with a structured safety notice; the
original response is preserved for audit purposes. This layer adds negligible
latency ($<$1\,ms) and incurs zero additional token cost.

\end{document}
